\documentclass[11pt,a4paper]{article}
\usepackage[margin=2.1cm]{geometry}
\usepackage[T1]{fontenc}
\usepackage{lmodern}
\usepackage{amsmath,amssymb,booktabs,graphicx,xcolor}
\usepackage{pgfplots}
\usepgfplotslibrary{groupplots,fillbetween}
\pgfplotsset{compat=1.18}
\usepackage[hidelinks]{hyperref}
\definecolor{MICEcolor}{HTML}{D95F59}
\definecolor{IP3Ocolor}{HTML}{2878B5}
\definecolor{CSPOcolor}{HTML}{38905B}
\setlength{\parindent}{0pt}
\setlength{\parskip}{0.5em}
\title{Goal 3: Recent Safe-RL Comparisons\\\large IP3O and CSPO in the MICE evaluation setting}
\author{RL Reproducibility Project}
\date{\today}
\begin{document}
\maketitle
\input{status.tex}
\section{Question and protocol}
The intended question is whether MICE's safety--reward benefits remain competitive
against recent online constrained reinforcement learning methods.
IP3O uses a smooth incentive--penalty objective; CSPO adds a correction informed
by local constraint sensitivity. The prepared algorithms and reference paper are
listed in the bibliography.

The configured experiment runs PointGoal1, CarGoal1, PointCircle1 and CarCircle1 with seeds
$\{0,1,2,3,4,5\}$ and an episode cost limit of 25. Each run collects seven million interactions. MICE and CSPO run 350 epochs of 20,000 steps. IP3O runs 1,400 epochs of 5,000 steps. Plots use environment steps, and tables
truncate each task to the smallest observed completed-episode horizon across
all included runs. No seed is extrapolated beyond its observed data.
This is a reduced-budget extension, not a full reproduction of MICE's
approximately ten-million-step experiments.

Algorithm-specific normalization, optimizer settings and update frequencies are
retained. IP3O's missing configuration parameter is supplied as $\alpha=0.1$
for navigation; CSPO uses its released $\alpha=0.3$.
All runs use vector observations and the same pinned environment versions.
Check \texttt{resolved\_config.json} in each run for the actual configuration;
command-line overrides may differ from these planned defaults.

\section{Metrics and interpretation}
For episode $i$ with length $L_i$, let
\[
R_i=\sum_{t=0}^{L_i-1}r_{i,t},\qquad C_i=\sum_{t=0}^{L_i-1}c_{i,t}.
\]
The plotted costs are extrinsic, raw environment costs; MICE's intrinsic cost
is not included. These sums are undiscounted. This explicit logging definition
should be distinguished from discounted-value terminology in the reference paper.
For $T$ training interactions and $N(T)$ completed episodes, report
\begin{align*}
A(T)&=\sum_{t=1}^{T} c_t &&\text{(accumulated training cost)},\\
V(T)&=\sum_{i=1}^{N(T)}\mathbf{1}\{C_i>d\} &&\text{(violating episode count)},\\
F(T)&=V(T)/N(T) &&\text{(violation fraction)},\\
E(T)&=\sum_{i=1}^{N(T)}\max(0,C_i-d) &&\text{(excess episode cost)}.
\end{align*}
$A(T)$ includes every recorded training step, including unfinished episode fragments.
The episode metrics include completed episodes only. Exceeding an episode limit
is a descriptive safety metric; an expected-cost constraint does not imply an
individual-episode guarantee.

Return and episode cost are smoothed within each seed before averaging seeds.
Bands show one population standard deviation across seeds, following the visual
convention of MICE Figure 3. They are not confidence intervals. No variability
band is drawn for a single seed. Cumulative metrics are not smoothed.
Interpolation is restricted to the common observed range of seeds for each
algorithm and task. The data tables and analysis notes specify the exact export.

\section{Results at a common interaction budget}
\input{results.tex}
\IfFileExists{timing.tex}{\input{timing.tex}}{}
\input{generated_figures.tex}
\clearpage
\input{discussion.tex}
\section{Reproducibility and limitations}
The run directory contains complete configurations, source revisions, package
versions, console output, per-episode and per-epoch CSVs, cumulative cost logs,
TensorBoard events and checkpoints. Files are flushed during training.
Interrupted runs retain partial observations and are identified in the data status.
The exported \texttt{included\_runs.json} records exactly which runs were analyzed.

Training return is not held-out final-policy evaluation. Run a separate common
evaluation protocol before drawing final-policy conclusions. Reduced-budget
results may reflect learning speed rather than asymptotic performance.
The study does not independently establish MICE's proposed causal explanation
about cost underestimation. Compatible MICE results must be included before
claiming that either comparator outperforms MICE.

\begin{thebibliography}{9}
\bibitem{mice} Gao et al. \emph{Controlling Underestimation Bias in Constrained
Reinforcement Learning for Safe Exploration}. ICML, 2025.
\url{https://github.com/ShiqingGao/MICE}
\bibitem{ip3o} Hazra, Dasgupta and Dey. \emph{Incentivizing Safer Actions in Policy
Optimization for Constrained Reinforcement Learning}. IJCAI, 2025.
\url{https://ijcai.org/proceedings/2025/592}
\bibitem{cspo} Belouadah, Kubler and Le Traon. \emph{CSPO: Constraint-Sensitive
Policy Optimization for Safe Reinforcement Learning}. ICML, 2026.
\url{https://proceedings.mlr.press/v306/belouadah26a.html}
\end{thebibliography}
\end{document}
